\chapter{Body}
\label{Chapter:Body}

This section demonstrates some of the features of this template, including:
\begin{enumerate}
    \item Equations;
    \item Short-cuts;
    \item Notes; and
    \item Objects;
\end{enumerate}

These will allow you to:
\begin{itemize}
    \item Automatically refer to numbered objects without having to manually cross-reference; and
    \item Save keystrokes for commonly used sequences;
\end{itemize}

The template uses the \verb|inline| option on the \verb|enumitem| package which allows inline lists. This can be useful to:
\begin{enumerate*}
    \item save space in the document;
    \item improve the flow of the text; and
    \item prevent run-on sentences;
\end{enumerate*}

\section{Equations}
By default, \LaTeX\ has equations which can be referred to using the \verb|\label| and \verb|\activeref| macros. For example, \ref{eqn:error}.

\begin{equation}\label{eqn:error}
    e(t) = PV(t) - SP(t)
\end{equation}

It is common to have both equations and reactions in your work. Thesis.tex provides functionality to distinguish reactions and mathematical equations. See \ref{rxn:delayed}.

\begin{reaction}\label{rxn:delayed}
    {^{87}Br} \underset{56 sec}{\stackrel{\beta^-}{\longrightarrow}} {^{87}Kr^{*}} \to {^{86}Kr + n}
\end{reaction}

If you would prefer to not differentiate equations and reactions, you can remove that functionality by removing the lines using \verb|\renewcommand{\theequation}| and \verb|\renewcommand{\thechemequation}|.

\section{Short-cuts}
\subsection{Acronyms}
You will likely use acronyms in your thesis. Define your acronyms in uidaho.cls using \verb|\newacronym{tag}{abbreviation}{full-text}|. The first time you use an acronym, use \verb|\acf{tag}| to render the full acronym, as such: \acf{npp}. Afterwards, you can use \verb|\acs{tag}| to only render the abbreviation (\acs{npp}), or \verb|\acl{tag}| to only render the full text (\acl{npp}). For the plural cases, use \verb|\acsp{tag}| or \verb|\aclp{tag}|. Any acronyms that you use will automatically be defined in the List of Acronyms. 

\subsection{Custom Commands}
If you have repeated text objects that are tedious to type in, \eg $^{135}Xe$, you can use the \verb|\newcommand| command to define a short-cut. uidaho.cls contains quite a few nuclides, as well as some latin phrases, \eg, \ie, \etal, and even scientific notation.

Generally, nuclides take one optional argument. For example, \verb|\Xe| renders \Xe, but \verb|\Xe[136]| renders \Xe[136]. \verb|6.02\sci[23]| will render 6.02\sci[23].

\section{Notes}
There are two types of notes supported by this template. Footnotes\footnotemark and margin notes\note{margin notes remind you to come back and do something}. Margin notes are set-up in Thesis.tex using \verb|\newcommandx|. You can customize the appearance of the notes and have different types of notes if you wish.

\footnotetext{Footnotes provide additional context to the reader.}


\section{Objects}
By use of the `cleveref' package, you don't need to type out `Figure' or `Table' when referring to your objects. It automatically adds the label if you use \verb|\cref| instead of \verb|\ref| (see \cref{fig:picture}).

\begin{figure}[ht!]
    \centering
    \includegraphics[width=0.75\textwidth]{ANS}
    \caption[Short version of caption]{Long version of caption. The short version, which is input in square brackets, is displayed in the List of Figures. LaTeX knows to look for graphics in the ./img/ directory because it is specified using the graphicspath command in Thesis.tex.}
    \label{fig:picture}
\end{figure}

This template also has support for tikz drawings (\cref{fig:tikz}) and pgf plots (\cref{fig:pgf}).

\begin{figure}[!ht]
    \centering
    
\begin{tikzpicture}
    
    %Sum
    \draw[->] (-3,0) node[anchor=east]{$SP(s)$}-- (-2,0);
    \draw (-1.75,0) circle (0.25)node{\scriptsize$-$};
    %Controller
    \draw[->] (-1.5,0) -- (-0.5,0)node[pos=0.5,anchor=south]{$e(s)$};
    \draw (-0.5,-0.5) rectangle (0.5,0.5) node[pos=0.5]{$C(s)$};
    \draw[->] (0.5,0) -- (1.5,0) node[pos=0.5,anchor=south]{$u(s)$};
    %Actuator
    \draw (1.5,-0.5) rectangle (2.5,0.5) node[pos=0.5]{$A(s)$};
    \draw[->] (2.5,0) -- (3.5,0);
    %Process
    \draw (3.5,-0.5) rectangle (4.5,0.5) node[pos=0.5]{$P(s)$};
    \draw[->] (4.5,0) -- (6.5,0) node[anchor=west]{$PV(s)$};
    %Transducer
    \draw[->] (5.5,0) -- (5.5,-1.5) -- (2.5,-1.5);
    \draw (1.5,-2) rectangle (2.5,-1) node[pos=0.5]{$H(s)$} ;
    \draw[->] (1.5,-1.5) -- (-1.75,-1.5) -- (-1.75,-0.25);
\end{tikzpicture}

    \caption[Feedback control loop]{Feedback control loop. The process-variable ($PV$) is measured by the transducer ($H$) and compared to the set-point ($SP$). The controller ($C$) uses the actuator ($A$) to control the process ($P$) based on the error ($e$).}
    \label{fig:tikz}
\end{figure}

\begin{figure}[!ht]
    \centering
    \subfloat[\centering Step-function]{{\resizebox{0.4\textwidth}{!}{% This file was created by matlab2tikz.
%
%The latest updates can be retrieved from
%  http://www.mathworks.com/matlabcentral/fileexchange/22022-matlab2tikz-matlab2tikz
%where you can also make suggestions and rate matlab2tikz.
%
\definecolor{mycolor1}{rgb}{0.00000,0.44700,0.74100}%
\definecolor{mycolor2}{rgb}{0.85000,0.32500,0.09800}%
%
\begin{tikzpicture}

\begin{axis}[axis lines = left,
width=4.521in,
height=3.566in,
at={(0.758in,0.481in)},
scale only axis,
xmin=0,
xmax=10,
xtick={ 0,  1,  2,  3,  4,  5,  6,  7,  8,  9, 10},
xticklabel style={font=\sansmath\sffamily},
xlabel style={font=\sansmath\sffamily},
xlabel={Time Constants, $\tau_F$},
ymin=-0.05,
ymax=1.05,
ylabel style={font=\sansmath\sffamily},
ylabel={Normalized Magnitude},
yticklabel style={font=\sansmath\sffamily},
ytick={0,0.632,0.865,0.95,1},
axis background/.style={fill=white},
legend style={at={(0.711,0.171)}, anchor=south west, legend cell align=left, align=left, font=\sansmath\sffamily}
]

\addplot [color=mycolor1]
  table[row sep=crcr]{%
0	0\\
0.2	0\\
0.4	0\\
0.6	0\\
0.8	0\\
0.999999999999993	0\\
1	1\\
1.00000000000001	1\\
1.20000000000001	1\\
1.40000000000001	1\\
1.60000000000001	1\\
1.80000000000001	1\\
2.00000000000001	1\\
2.20000000000001	1\\
2.40000000000001	1\\
2.60000000000001	1\\
2.80000000000001	1\\
3.00000000000002	1\\
3.20000000000002	1\\
3.40000000000002	1\\
3.60000000000002	1\\
3.80000000000002	1\\
4.00000000000002	1\\
4.20000000000002	1\\
4.40000000000002	1\\
4.60000000000002	1\\
4.80000000000002	1\\
4.99999999999996	1\\
5.00000000000002	1\\
5.20000000000002	1\\
5.40000000000002	1\\
5.60000000000002	1\\
5.80000000000002	1\\
6.00000000000002	1\\
6.20000000000002	1\\
6.40000000000002	1\\
6.60000000000002	1\\
6.80000000000002	1\\
7.00000000000002	1\\
7.20000000000002	1\\
7.40000000000002	1\\
7.60000000000002	1\\
7.80000000000002	1\\
8.00000000000002	1\\
8.20000000000002	1\\
8.40000000000002	1\\
8.60000000000002	1\\
8.80000000000002	1\\
9.00000000000002	1\\
9.20000000000002	1\\
9.40000000000001	1\\
9.60000000000001	1\\
9.80000000000001	1\\
10	1\\
};
\addlegendentry{Input}

\addplot [color=mycolor2, dashed]
  table[row sep=crcr]{%
0	0\\
0.2	0\\
0.4	0\\
0.6	0\\
0.8	0\\
0.999999999999993	0\\
1	0\\
1.00000000000001	1.42108547152019e-14\\
1.20000000000001	0.181269226666678\\
1.40000000000001	0.329679920797011\\
1.60000000000001	0.451188323173276\\
1.80000000000001	0.550670991417293\\
2.00000000000001	0.63212051332198\\
2.20000000000001	0.698805743378635\\
2.40000000000001	0.753402993352831\\
2.60000000000001	0.798103442046079\\
2.80000000000001	0.834701074973048\\
3.00000000000002	0.864664683281515\\
3.20000000000002	0.889196811483763\\
3.40000000000002	0.909282019778302\\
3.60000000000002	0.925726397897851\\
3.80000000000002	0.939189916312655\\
4.00000000000002	0.950212913156196\\
4.20000000000002	0.959237779886358\\
4.40000000000002	0.966626716003575\\
4.60000000000002	0.972676265384934\\
4.80000000000002	0.977629217628252\\
4.99999999999996	0.981684352048706\\
5.00000000000002	0.981684352048707\\
5.20000000000002	0.985004415388737\\
5.40000000000002	0.987722653414635\\
5.60000000000002	0.989948158535683\\
5.80000000000002	0.991770248064495\\
6.00000000000002	0.993262048833503\\
6.20000000000002	0.994483432030772\\
6.40000000000002	0.995483416040408\\
6.60000000000002	0.996302133721938\\
6.80000000000002	0.996972443082479\\
7.00000000000002	0.997521245983607\\
7.20000000000002	0.997970567807256\\
7.40000000000002	0.998338441411407\\
7.60000000000002	0.998639630851823\\
7.80000000000002	0.998886223915294\\
8.00000000000002	0.999088117244848\\
8.20000000000002	0.999253413526685\\
8.40000000000002	0.999388746679343\\
8.60000000000002	0.999499548096076\\
8.80000000000002	0.999590264625684\\
9.00000000000002	0.999664537040124\\
9.20000000000002	0.999725346151436\\
9.40000000000001	0.999775132442167\\
9.60000000000001	0.999815894010477\\
9.80000000000001	0.999849266760823\\
10	0.999876590058521\\
};
\addlegendentry{Output}

\end{axis}


\end{tikzpicture}%}}}
    \qquad
    \subfloat[\centering Ramp-function]{\resizebox{0.4\textwidth}{!}{% This file was created by matlab2tikz.
%
%The latest updates can be retrieved from
%  http://www.mathworks.com/matlabcentral/fileexchange/22022-matlab2tikz-matlab2tikz
%where you can also make suggestions and rate matlab2tikz.
%
\definecolor{mycolor1}{rgb}{0.00000,0.44700,0.74100}%
\definecolor{mycolor2}{rgb}{0.85000,0.32500,0.09800}%
%
\begin{tikzpicture}

  \begin{axis}[axis lines = left,
    width=4.521in,
    height=3.566in,
    at={(0.758in,0.481in)},
    scale only axis,
    xmin=0,
    xmax=10,
    xtick={ 0,  1,  2,  3,  4,  5,  6,  7,  8,  9, 10},
    xticklabel style={font=\sansmath\sffamily},
    xlabel style={font=\sansmath\sffamily},
    xlabel={Time Constants, $\tau_F$},
    ymin=-0.05,
    ymax=1.05,
    ylabel style={font=\sansmath\sffamily},
    ylabel={Normalized Magnitude},
    yticklabel style={font=\sansmath\sffamily},
    ytick={0,0.632,0.865,0.95,1},
    axis background/.style={fill=white},
    legend style={at={(0.711,0.171)}, anchor=south west, legend cell align=left, align=left, font=\sansmath\sffamily}
    ]
    
\addplot [color=mycolor1]
  table[row sep=crcr]{%
0	0\\
0.2	0\\
0.4	0\\
0.6	0\\
0.8	0\\
0.999999999999993	0\\
1	0\\
1.00000000000001	3.5527136788005e-15\\
1.20000000000001	0.0500000000000035\\
1.40000000000001	0.100000000000004\\
1.60000000000001	0.150000000000004\\
1.80000000000001	0.200000000000004\\
2.00000000000001	0.250000000000004\\
2.20000000000001	0.300000000000004\\
2.40000000000001	0.350000000000004\\
2.60000000000001	0.400000000000004\\
2.80000000000001	0.450000000000004\\
3.00000000000002	0.500000000000004\\
3.20000000000002	0.550000000000004\\
3.40000000000002	0.600000000000004\\
3.60000000000002	0.650000000000004\\
3.80000000000002	0.700000000000004\\
4.00000000000002	0.750000000000004\\
4.20000000000002	0.800000000000004\\
4.40000000000002	0.850000000000004\\
4.60000000000002	0.900000000000004\\
4.80000000000002	0.950000000000004\\
4.99999999999996	0.99999999999999\\
5.00000000000002	1\\
5.20000000000002	1\\
5.40000000000002	1\\
5.60000000000002	1\\
5.80000000000002	1\\
6.00000000000002	1\\
6.20000000000002	1\\
6.40000000000002	1\\
6.60000000000002	1\\
6.80000000000002	1\\
7.00000000000002	1\\
7.20000000000002	1\\
7.40000000000002	1\\
7.60000000000002	1\\
7.80000000000002	1\\
8.00000000000002	1\\
8.20000000000002	1\\
8.40000000000002	1\\
8.60000000000002	1\\
8.80000000000002	1\\
9.00000000000002	1\\
9.20000000000002	1\\
9.40000000000001	1\\
9.60000000000001	1\\
9.80000000000001	1\\
10	1\\
};
\addlegendentry{Input}

\addplot [color=mycolor2, dashed]
  table[row sep=crcr]{%
0	0\\
0.2	0\\
0.4	0\\
0.6	0\\
0.8	0\\
0.999999999999993	0\\
1	0\\
1.00000000000001	2.52378540224252e-29\\
1.20000000000001	0.00468269333333398\\
1.40000000000001	0.0175800198007507\\
1.60000000000001	0.0372029192066846\\
1.80000000000001	0.0623322521456803\\
2.00000000000001	0.0919698716695085\\
2.20000000000001	0.125298564155345\\
2.40000000000001	0.161649251661796\\
2.60000000000001	0.200474139488484\\
2.80000000000001	0.241324731256742\\
3.00000000000002	0.283833829179625\\
3.20000000000002	0.327700797129063\\
3.40000000000002	0.372679495055428\\
3.60000000000002	0.418568400525541\\
3.80000000000002	0.46520252092184\\
4.00000000000002	0.512446771710955\\
4.20000000000002	0.560190555028414\\
4.40000000000002	0.60834332099911\\
4.60000000000002	0.656830933653771\\
4.80000000000002	0.705592695592941\\
4.99999999999996	0.754578911987813\\
5.00000000000002	0.754578911987827\\
5.20000000000002	0.799066202819486\\
5.40000000000002	0.835489316845594\\
5.60000000000002	0.865310041159399\\
5.80000000000002	0.8897251858382\\
6.00000000000002	0.90971461612212\\
6.20000000000002	0.926080577836966\\
6.40000000000002	0.939479894328106\\
6.60000000000002	0.950450327081035\\
6.80000000000002	0.959432157972642\\
7.00000000000002	0.966785859324477\\
7.20000000000002	0.972806560919127\\
7.40000000000002	0.977735894591724\\
7.60000000000002	0.981771691761507\\
7.80000000000002	0.98507592309934\\
8.00000000000002	0.987781198977837\\
8.20000000000002	0.989996091589918\\
8.40000000000002	0.991809492331058\\
8.60000000000002	0.993294179322215\\
8.80000000000002	0.994509738250642\\
9.00000000000002	0.995504953752146\\
9.20000000000002	0.996319767309325\\
9.40000000000001	0.996986880243117\\
9.60000000000001	0.997533066131301\\
9.80000000000001	0.997980245325918\\
10	0.998346364693745\\
};
\addlegendentry{Output}

\end{axis}


\end{tikzpicture}%}}
    \caption[Pre-filter on (a) step-function and (b) ramp-function]{Pre-filter on step-function and ramp-function. When the pre-filter acts on a step-function, it follows an exponential curve, reaching 63.2\% of the magnitude of the step in 1 time constant, 86.5\% in 2 time constants, 95.0\% in 3 time constants and so on. The ramp-function exhibits similar but more complicated dynamics due to the changing input.}
    \label{fig:pgf}
\end{figure}

Tables require you to use \verb|\Cref| instead of \verb|\cref| (see \Cref{tab:params})

\begin{table}[ht!]
    \caption[Relevant nuclear constants]{Relevant nuclear constants \cite{Lamarsh}. The \I fission yield ($\gamma$) is the sum of direct and indirect fission. The \I microscopic neutron capture cross-section ($\sigma$) is neglected as it is so small that it is insignificant compared to its own decay rate and the \Xe cross-section.}
    \centering\begin{tabular}{c|ccc}
                   &  $\gamma \;(\%)$ &  $\lambda \; (hr^{-1})$ &  $\sigma_a \; (Mb)$ \\ \hline
        \I  & 6.39            & 0.1035                 & -                \\
        \Xe & 0.237           & 0.0753                 & 2.65\textsuperscript{*}
    \end{tabular}\\
    \textsuperscript{*}At 0.025 eV
    \label{tab:params}
\end{table}

The template also has support for Python and Serpent code blocks using the custom `./rcs/slither.sty' package. Codes are included in Appendix \ref{app:codes}.

Finally, LaTeX handles citations for you. Include the data about your citations in ./rcs/References.bib, and cite them here using \verb|\cite|. For example, \cite{RootMS}. If you are citing textbooks and want to include a chapter, put any such information in square brackets \cite[Ch. 6]{Lamarsh}.